\section{地图与定义：Base、Instruction、Preference 三层不要混淆}

上一章说明为什么出现对齐，本章说明讨论对象。开放社区常把 ``fine-tuned''、``chat''、``instruct''、``aligned'' 和 ``RLHF model'' 互换使用；课程则要求区分训练数据、损失函数与模型接口。只有先分清层次，后面的模型比较和评测才有意义。

\subsection{课程 atlas 与四个章节}

本节先看讲者提供的模型集合和章节 atlas。模型图标不是完整谱系，而是一条可复现线索：读者可沿 Hugging Face collection 找到许多课堂模型；四章则按 instruction、evaluation、RLHF 与 modern ecosystem 展开，暗示算法之外还有数据与测量基础设施。

\lecturefigure{slide-017.jpg}{课程 atlas：模型集合、参考资源与发布时间线}{official deck, p.~17}

\lecturefigure{slide-019.jpg}{第一章焦点：instruction models}{official deck, p.~19}

\lecturefigure{slide-020.jpg}{第二章加入：expectations 与 evaluations}{official deck, p.~20}

\lecturefigure{slide-021.jpg}{第三章加入：RLHF 与 preference optimization}{official deck, p.~21}

\lecturefigure{slide-022.jpg}{第四章加入：modern ecosystem}{official deck, p.~22}

\lecturefigure{slide-023.jpg}{完整章节地图：instruction、evals、RLHF、ecosystem}{official deck, p.~23}

\begin{knowledgebox}{读图：为什么 evaluation 被单列一章}
没有快速评测，开放团队无法知道一次数据或优化变化是否值得发布；没有长期人类偏好信号，快速自动评测又会被投机。Evaluation 不是训练后的附属步骤，而是决定团队能否迭代、哪些模型被看见、哪些数据被继续收集的反馈通道。
\end{knowledgebox}

\subsection{Base model 与 aligned model}

现在把模型分成两层。Base model 主要通过大规模 next-token prediction 学习通用分布；经过 alignment、fine-tuning 或 preference training 的模型在其上继续训练，使输出适配指令格式、对话角色和偏好。后者可以在较少计算下改变交互体验，但能力上限仍强烈受前者约束。

\lecturefigure{slide-024.jpg}{Base models 是开放对齐生态的能力底座}{official deck, p.~24}

\lecturefigure{slide-025.jpg}{Aligned、fine-tuned 与 preference-trained models 是行为适配层}{official deck, p.~25}

\begin{importantbox}{Base 与 aligned 的条件概率差异}
Base model 学习 $p_{\theta_0}(x_t\mid x_{<t})$；对齐阶段得到参数 $\theta$，使目标交互分布下的回答概率重新排序。可概念性写成
\[
p_{\theta}(y\mid x)\propto p_{\theta_0}(y\mid x)\exp\!\left(\frac{r(x,y)}{\beta}\right),
\]
其中 $r(x,y)$ 是对回答的偏好或奖励，$\beta$ 控制偏离 base/reference distribution 的程度。公式表达“保留能力同时偏向更受欢迎回答”的目标，不代表所有方法都显式这样实现。
\end{importantbox}

\teachervoice{00:08:48--00:09:38，讲者主动承认时间线省略了大量学术与基础设施贡献，并强调没有 LLaMA/Llama 2 等 base model 就没有后续开放对齐生态。Alignment model 不是脱离底座独立存在的产品。}

\subsection{对齐术语表}

本节把常见术语按“数据—目标—输出”拆开。Instruction fine-tuning 与 supervised fine-tuning 常使用同一 autoregressive loss，但前者强调指令接口；RLHF 使用人类偏好学习奖励并优化 policy；DPO 直接从偏好对优化 policy/reference log-ratio。Alignment 则是更宽泛的目标概念。

\lecturefigure{slide-026.jpg}{Instruction fine-tuning、SFT、RLHF、DPO 与 alignment 的课堂定义}{official deck, p.~26}

\begin{knowledgebox}{术语消化：名称、数据与机制}
\begin{tabularx}{\textwidth}{L{0.18\textwidth} L{0.24\textwidth} X}
\toprule
术语 & 主要数据 & 核心机制 \\
\midrule
IFT & instruction-response pairs & 用 LM loss 学会指令格式、system prompt、多轮角色与目标风格。 \\
SFT & 任意带标签示范 & 用监督损失模仿目标输出；不一定是自然语言指令。 \\
RLHF & preference comparisons + prompts & 学 reward model，再用 PPO 等方法最大化奖励并限制 policy drift。 \\
DPO & chosen/rejected response pairs & 直接优化偏好 log-ratio，不显式运行 reward-model rollout loop。 \\
Alignment & 用户/组织期望的行为标准 & 可由多种数据、损失、规则和系统机制共同实现。 \\
\bottomrule
\end{tabularx}
\end{knowledgebox}

\teachervoice{00:09:53--00:11:14，Nathan 把 alignment 定义得很宽：任何让模型更贴近用户 desires 的训练都可纳入，而 IFT、SFT、RLHF 与 DPO 只是不同机制。讲义沿用这个宽定义，但要求每次具体说明数据和损失。}

\subsection{Chapter 0：复现 ChatGPT 的抢跑期}

有了术语，课程回到 2023 年初。当时的开放社区并不知道 dialogue model 的关键配方，也缺少成熟 red-teaming、chat data 与评测。大量团队快速组合 base model、合成数据和轻量 fine-tuning，形成“每周一个模型”的抢跑期。

\lecturefigure{slide-027.jpg}{Chapter 0：复现 ChatGPT 的 land grab 与基础问题}{official deck, p.~27}

\begin{warningbox}{不要用后见之明评价早期模型}
今天看起来基本的 chat template、system prompt、拒答数据和 pairwise eval，在当时都没有统一实现。早期结果的价值常在于暴露接口和数据问题，而不是最终分数。历史分析应问“它解锁了什么”，而不只是“它现在是否过时”。
\end{warningbox}

\teachervoice{00:11:30--00:12:30，讲者把这段时期称为 land grab：大家一边发布模型，一边才弄清什么是 dialogue、red-teaming 和有用回答。这个背景解释了后续数据与评测为何迅速成为主战场。}

\subsection{本章小结}

本章建立了后续分析需要的分层：base model 决定能力底座，IFT/SFT 改变接口与模仿行为，RLHF/DPO 使用偏好重新排序回答，alignment 是更宽的系统目标。课程 atlas 则把数据、评测和算法并列，提醒我们不要把模型进展归因于单个优化器。

